\section{Behavior Under Algorithmic Redistricting}
\label{sec:behavior}

\subsection{Why Compactness Reduces Declination}

Declination is elevated when Republican-won districts are clustered near 50\%
(cracking: $\bar{v}_R^-$ close to 50\%, making $\theta_R$ small) while Democratic-won
districts are clustered far from 50\% (packing: $\bar{v}_D^+$ far above 50\%,
making $\theta_D$ large). The angular difference $\theta_D - \theta_R$ is large
positive when cracking and packing are both present simultaneously.

Cracking requires drawing Republican-favorable districts at exactly 52--55\% Democratic
vote share (just above Democrat-losing), which requires careful geographic engineering:
the district must include enough Democratic voters to look competitive without including
so many that it flips. This level of precision requires access to fine-grained partisan
data at the precinct or block level---the mapmaker must select exactly the right
precincts to push the district to the desired vote share.

\textsc{Bisect} lacks this capability by design. The algorithm receives only population
counts and geographic adjacency data; it cannot target specific vote-share outcomes.
Districts are defined by recursive geographic bisection: the algorithm cuts the state
in half geographically, then cuts each half in half, and so on until individual districts
are formed. The resulting districts have whatever partisan composition results from the
geographic distribution of voters within their boundaries---which, for compact districts
in states without extreme partisan segregation, tends to be a mixture that does not
cluster at any particular vote share.

The result is that \textsc{bisect} maps produce approximately symmetric distributions
of Republican-won and Democratic-won district vote shares---both parties' won districts
are similarly close to 50\% on average, producing $\theta_D \approx \theta_R$ and
$\delta \approx 0$. The residual $|\delta| \approx 0.08$--$0.10$ reflects the
geographic sorting of NC's urban Democratic cores into high-margin Democratic districts:
those districts are far from 50\% ($\bar{v}_D^+ \approx 0.70$), pulling $\theta_D$
above $\theta_R$ even in the absence of deliberate cracking.

\subsection{Declination as the Most Sensitive Metric}

Among the six L-series metrics, Declination shows the largest percentage reduction
from enacted to \textsc{bisect} maps (70--80\%). This reflects that Declination is
most directly sensitive to the geometric signature of deliberate cracking---the
clustering of Republican-won districts near 50\%---which compactness most directly
prevents. EG and MM are aggregates of the full vote-share distribution; Partisan Bias
measures only the outcome at the 50\% pivot; Declination measures the relative
positions of both parties' won districts simultaneously, making it uniquely sensitive
to the cracking pattern.

The practical implication is that Declination provides the strongest discrimination
between compact neutral maps and deliberately cracked enacted maps---but at the cost
of the three methodological limitations documented in Section~\ref{sec:definition}.
The high discrimination power and the methodological fragility are two sides of the
same coin: Declination's sensitivity to the cracking-specific geometric pattern means
it is also sensitive to other geometric features (boundary effects, small district
counts, near-50\% vote shares) that produce false positives.
